\documentclass[11pt]{article}
\usepackage[utf8]{inputenc}
\usepackage[T1]{fontenc}
\usepackage[french]{babel}
\usepackage{geometry}
\usepackage{array}
\geometry{margin=1in}

\begin{document}
\selectlanguage{french}

\begin{center}
  {\LARGE \textbf{Vocabulaire : la salle de classe}}\\[6pt]
  {\large Singulier / Pluriel — Un / Une / Des}\\[12pt]
  \small (Les mots marqués en \textbf{gras} sont ceux indiqués comme \emph{starred} dans la liste.)
\end{center}

\vspace{8pt}

% --- Tableau pour le masculin ---
\begin{center}
\begin{tabular}{|p{7.5cm}|p{7.5cm}|}
\hline
\multicolumn{2}{|c|}{\textbf{AU MASCULIN}}\\
\hline
\textbf{SINGULIER} & \textbf{PLURIEL} \\
\hline
Un crayon & Des crayons \\
\hline
Un élève & Des élèves \\
\hline
Un livre & Des livres \\
\hline
Un sac à dos & Des sacs à dos \\
\hline
Un classeur & Des classeurs \\
\hline
Un tableau (de photos / d'affichage) & Des tableaux \\
\hline
Un meuble de rangement & Des meubles de rangement \\
\hline
Un plafond & Des plafonds \\
\hline
Un téléphone portable / un cellulaire & Des téléphones portables / des cellulaires \\
\hline
Un ordinateur & Des ordinateurs \\
\hline
Un crayon de couleur & Des crayons de couleur \\
\hline
Un crayon gras & Des crayons gras \\
\hline
Un bureau (office desk) & Des bureaux \\
\hline
Un pupitre (desk — élève) & Des pupitres \\
\hline
Un dictionnaire & Des dictionnaires \\
\hline
Un drapeau (français / américain) & Des drapeaux \\
\hline
Un sol & Des sols \\
\hline
Un dossier (folder) & Des dossiers \\
\hline
Un casque audio & Des casques audio \\
\hline
Un feutre / un marqueur & Des feutres / des marqueurs \\
\hline
Un cahier & Des cahiers \\
\hline
Un papier (ou \emph{une feuille de papier}) & Des papiers (ou des feuilles de papier) \\
\hline
Un trombone & Des trombones \\
\hline
Un stylo & Des stylos \\
\hline
Un taille-crayon & Des taille-crayons \\
\hline
Un écran & Des écrans \\
\hline
Un ruban (adhésif) & Des rubans \\
\hline
Un mur & Des murs \\
\hline
\end{tabular}
\end{center}

\vspace{10pt}

% --- Tableau pour le féminin ---
\begin{center}
\begin{tabular}{|p{7.5cm}|p{7.5cm}|}
\hline
\multicolumn{2}{|c|}{\textbf{AU FÉMININ}}\\
\hline
\textbf{SINGULIER} & \textbf{PLURIEL} \\
\hline
Une table & Des tables \\
\hline
Une élève & Des élèves \\
\hline
Une professeure / une prof & Des professeures / des profs \\
\hline
Une salle de classe & Des salles de classe \\
\hline
Une étagère & Des étagères \\
\hline
Une calculatrice & Des calculatrices \\
\hline
Une chaise & Des chaises \\
\hline
Une craie & Des craies \\
\hline
Une horloge & Des horloges \\
\hline
Une porte & Des portes \\
\hline
Une gomme & Des gommes \\
\hline
Une colle (ou un bâton de colle) & Des bâtons de colle \\
\hline
Une perforatrice & Des perforatrices \\
\hline
Une lampe & Des lampes \\
\hline
Une carte / un plan & Des cartes / des plans \\
\hline
Une trousse à crayons & Des trousses à crayons \\
\hline
Une affiche & Des affiches \\
\hline
Une imprimante & Des imprimantes \\
\hline
Une règle & Des règles \\
\hline
Une agrafeuse & Des agrafeuses \\
\hline
Une télévision (la télé) & Des télévisions \\
\hline
Une poubelle / une corbeille & Des poubelles / des corbeilles \\
\hline
Une fenêtre & Des fenêtres \\
\hline
\end{tabular}
\end{center}

\vspace{8pt}

\noindent\textbf{Cas particuliers / notes :}
\begin{itemize}
  \item \textbf{Les ciseaux} : on dit souvent \emph{une paire de ciseaux} (singulier concept) ou \emph{des ciseaux} (pluriel).
  \item \textbf{Les écouteurs} : usage courant au pluriel (\emph{des écouteurs}).
  \item \textbf{L'élève / L'étudiant(e)} : l'article change selon le genre — \emph{un élève} / \emph{une élève}.
  \item Correction de coquilles de la liste d'origine : \textit{bes liveres} $\to$ \textbf{des livres}, \textit{mouchairs} $\to$ \textbf{mouchoirs}.
\end{itemize}

\vspace{6pt}
\noindent\textbf{Vocabulaire à pratiquer (sélection) :}\\
\noindent \textbf{Masculin:} \textit{un crayon, un élève, un livre, un sac à dos, un classeur, un tableau, un ordinateur, un bureau, un pupitre, un stylo, un cahier, un taille-crayon.}\\
\noindent \textbf{Féminin:} \textit{une table, une élève, une étagère, une calculatrice, une chaise, une gomme, une trousse à crayons, une règle, une lampe, une poubelle, une fenêtre.}


\end{document}
